\chapter{Subversion Quick-Start Guide}\label{svn.intro}

If you\textquotesingle re eager to get Subversion up and running (and
you enjoy learning by experimentation), this appendix will show you how
to create a repository, import code, and then check it back out again as
a working copy. Along the way, we give links to the relevant chapters of
this book.

If you\textquotesingle re new to the entire concept of version control
or to the ``copy-modify-merge'' model used by both CVS and Subversion,
you should read \hyperref[svn.basic]{Fundamental Concepts} before going
any further.

\section{Installing Subversion}\label{svn.intro.install}

Subversion is built on a portability layer called APR---the Apache
Portable Runtime library. The APR library provides all the interfaces
that Subversion needs to function on different operating systems: disk
access, network access, memory management, and so on. The abstraction
layer provided by APR enables Subversion clients and servers to run on
any operating system that other APR-based applications run on: Windows,
Linux, all flavors of BSD, Mac OS X, NetWare, and others.

Although the APR library is part of the Apache HTTP Server (or,
\texttt{httpd}), and \texttt{httpd} can be configured to serve
Subversion repositories, \texttt{httpd} is \emph{not} a required
component of a Subversion installation.

The easiest way to get Subversion is to download a binary package built
for your operating system. Subversion\textquotesingle s web site
(\href{https://subversion.apache.org}{}) often has these packages
available for download, posted by volunteers. The site usually contains
graphical installer packages for users of Microsoft operating systems.
If you run a Unix-like operating system, you can use your
system\textquotesingle s native package distribution system (RPMs, DEBs,
the ports tree, etc.) to get Subversion.

Alternatively, you can build Subversion directly from source code,
though it\textquotesingle s not always an easy task. (If
you\textquotesingle re not experienced at building open source software
packages, you\textquotesingle re probably better off downloading a
binary distribution instead!) From the Subversion web site, download the
latest source code release. After unpacking it, follow the instructions
in the \texttt{INSTALL} file to build it.

If you\textquotesingle re one of those folks that likes to use
bleeding-edge software, you can also get the Subversion source code from
the Subversion repository in which it lives. Obviously,
you\textquotesingle ll need to already have a Subversion client on hand
to do this. But once you do, you can check out a working copy from
\href{https://svn.apache.org/repos/asf/subversion}{}\footnote{Note that
  the URL checked out in the example ends not with \texttt{subversion},
  but with a subdirectory thereof called \texttt{trunk}. See our
  discussion of Subversion\textquotesingle s branching and tagging model
  for the reasoning behind this.}:

\begin{verbatim}
$ svn checkout https://svn.apache.org/repos/asf/subversion/trunk subversion
A    subversion/HACKING
A    subversion/INSTALL
A    subversion/README
A    subversion/autogen.sh
A    subversion/build.conf
…
\end{verbatim}

The preceding command will create a working copy of the latest
(unreleased) Subversion source code into a subdirectory named
\texttt{subversion} in your current working directory. You can adjust
that last argument as you see fit. Regardless of what you call the new
working copy directory, though, after this operation completes, you will
now have the Subversion source code. Of course, you will still need to
fetch a few helper libraries (apr, apr-util, etc.)---see the
\texttt{INSTALL} file in the top level of the working copy for details.

\section{High-Speed Tutorial}\label{svn.intro.quickstart}

\begin{quote}
``Please make sure your seat backs are in their full, upright position
and that your tray tables are stored. Flight attendants, prepare for
take-off\ldots.''
\end{quote}

What follows is a quick tutorial that walks you through some basic
Subversion configuration and operation. When you finish it, you should
have a general understanding of Subversion\textquotesingle s typical
usage.

The examples used in this appendix assume that you have \texttt{svn},
the Subversion command-line client, and \texttt{svnadmin}, the
administrative tool, ready to go on a Unix-like operating system. (This
tutorial also works at the Windows command-line prompt, assuming you
make some obvious tweaks.) We also assume you are using Subversion 1.2
or later (run \texttt{svn\ -\/-version} to check).

Subversion stores all versioned data in a central repository. To begin,
create a new repository:

\begin{verbatim}
$ cd /var/svn
$ svnadmin create repos
$ ls repos
conf/  dav/  db/  format  hooks/  locks/  README.txt
$
\end{verbatim}

This command creates a Subversion repository in the directory
\texttt{/var/svn/repos}, creating the \texttt{repos} directory itself if
it doesn\textquotesingle t already exist. This directory contains (among
other things) a collection of database files. You won\textquotesingle t
see your versioned files if you peek inside. For more information about
repository creation and maintenance, see
\hyperref[svn.reposadmin]{Repository Administration}.

Subversion has no concept of a ``project''. The repository is just a
virtual versioned filesystem, a large tree that can hold anything you
wish. Some administrators prefer to store only one project in a
repository, and others prefer to store multiple projects in a repository
by placing them into separate directories. We discuss the merits of each
approach in \hyperref[svn.reposadmin.projects.chooselayout]{Planning
Your Repository Organization}. Either way, the repository manages only
files and directories, so it\textquotesingle s up to humans to interpret
particular directories as ``projects''. So while you might see
references to projects throughout this book, keep in mind that
we\textquotesingle re only ever talking about some directory (or
collection of directories) in the repository.

In this example, we assume you already have some sort of project (a
collection of files and directories) that you wish to import into your
newly created Subversion repository. Begin by organizing your data into
a single directory called \texttt{myproject} (or whatever you wish). For
reasons explained in \hyperref[svn.branchmerge]{Branching and Merging},
your project\textquotesingle s tree structure should contain three
top-level directories named \texttt{branches}, \texttt{tags}, and
\texttt{trunk}. The \texttt{trunk} directory should contain all of your
data, and the \texttt{branches} and \texttt{tags} directories should be
empty:

\begin{verbatim}
/tmp/
   myproject/
      branches/
      tags/
      trunk/
         foo.c
         bar.c
         Makefile
         …
\end{verbatim}

The \texttt{branches}, \texttt{tags}, and \texttt{trunk} subdirectories
aren\textquotesingle t actually required by Subversion.
They\textquotesingle re merely a popular convention that
you\textquotesingle ll most likely want to use later on.

Once you have your tree of data ready to go, import it into the
repository with the \texttt{svn\ import} command (see
\hyperref[svn.tour.importing]{Getting Data into Your Repository}):

\begin{verbatim}
$ svn import /tmp/myproject file:///var/svn/repos/myproject \
      -m "initial import"
Adding         /tmp/myproject/branches
Adding         /tmp/myproject/tags
Adding         /tmp/myproject/trunk
Adding         /tmp/myproject/trunk/foo.c
Adding         /tmp/myproject/trunk/bar.c
Adding         /tmp/myproject/trunk/Makefile
…
Committed revision 1.
$ 
\end{verbatim}

Now the repository contains this tree of data. As mentioned earlier, you
won\textquotesingle t see your files by directly peeking into the
repository; they\textquotesingle re all stored within a database. But
the repository\textquotesingle s imaginary filesystem now contains a
top-level directory named \texttt{myproject}, which in turn contains
your data.

Note that the original \texttt{/tmp/myproject} directory is unchanged;
Subversion is unaware of it. (In fact, you can even delete that
directory if you wish.) To start manipulating repository data, you need
to create a new ``working copy'' of the data, a sort of private
workspace. Ask Subversion to ``check out'' a working copy of the
\texttt{myproject/trunk} directory in the repository:

\begin{verbatim}
$ svn checkout file:///var/svn/repos/myproject/trunk myproject
A    myproject/foo.c
A    myproject/bar.c
A    myproject/Makefile
…
Checked out revision 1.
$
\end{verbatim}

Now you have a personal copy of part of the repository in a new
directory named \texttt{myproject}. You can edit the files in your
working copy and then commit those changes back into the repository.

\begin{itemize}
\item
  Enter your working copy and edit a file\textquotesingle s contents.
\item
  Run \texttt{svn\ diff} to see unified diff output of your changes.
\item
  Run \texttt{svn\ commit} to commit the new version of your file to the
  repository.
\item
  Run \texttt{svn\ update} to bring your working copy ``up to date''
  with the repository.
\end{itemize}

For a full tour of all the things you can do with your working copy,
read \hyperref[svn.tour]{Basic Usage}.

At this point, you have the option of making your repository available
to others over a network. See \hyperref[svn.serverconfig]{Server
Configuration} to learn about the different sorts of server processes
available and how to configure them.

